\section{Theta Functions and Modular Forms}

\subsection{Classical Theta Functions}

The four Jacobi theta functions form the basis for all elliptic constructions:

\begin{definition}[Jacobi Theta Functions]
\index{theta functions!Jacobi|textbf}
\begin{align}
\vartheta_{00}(z|\tau) &\equiv \vartheta_3(z|\tau) = \sum_{n \in \mathbb{Z}} q^{n^2/2} e^{2\pi inz} \\
\vartheta_{01}(z|\tau) &\equiv \vartheta_4(z|\tau) = \sum_{n \in \mathbb{Z}}(-1)^n q^{n^2/2} e^{2\pi inz} \\
\vartheta_{10}(z|\tau) &\equiv \vartheta_2(z|\tau) = \sum_{n \in \mathbb{Z}} q^{(n+1/2)^2/2} e^{2\pi i(n+1/2)z} \\
\vartheta_{11}(z|\tau) &\equiv \vartheta_1(z|\tau) = -i\sum_{n \in \mathbb{Z}}(-1)^n q^{(n+1/2)^2/2} e^{2\pi i(n+1/2)z}
\end{align}
where $q = e^{2\pi i\tau}$ is the nome.
\end{definition}

\subsection{Modular Transformation Laws}

Under the generators of $SL_2(\mathbb{Z})$:
$$T: \tau \mapsto \tau + 1, \quad S: \tau \mapsto -1/\tau$$

The theta functions transform as:
\begin{align}
\vartheta_{ab}(z|\tau+1) &= e^{\pi ia^2/4} \vartheta_{a,\, b+1-a}(z|\tau) \\
\vartheta_{ab}(z/\tau|-1/\tau) &= (-i\tau)^{1/2} e^{\pi iz^2/\tau} \sum_{cd} K_{ab,cd} \vartheta_{cd}(z|\tau)
\end{align}
where the kernel matrix is $K_{ab,cd} = \frac{1}{2}(-1)^{ac+bd}e^{-\pi i(a+c)/2}$, giving explicitly:
$$K = \frac{1}{2}\begin{pmatrix} 1 & 1 & 1 & 1 \\ 1 & 1 & -1 & -1 \\ 1 & -1 & -i & i \\ 1 & -1 & i & -i \end{pmatrix}$$
in the basis $(\vartheta_{00}, \vartheta_{01}, \vartheta_{10}, \vartheta_{11})$.

\subsection{Higher Genus Theta Functions}

For genus $g$, theta functions depend on $g \times g$ period matrices $\Omega$:
$$\Theta[\epsilon](z|\Omega) = \sum_{n \in \mathbb{Z}^g} \exp\left[\pi i(n+\epsilon')^t\Omega(n+\epsilon') + 2\pi i(n+\epsilon')^t(z+\epsilon'')\right]$$
where $\epsilon = (\epsilon', \epsilon'') \in (\mathbb{Z}_2)^{2g}$ is the characteristic.

\subsection{Elliptic and Siegel Modular Forms}

\begin{definition}[Weight $k$ Modular Form]
\index{modular form|textbf}
\index{Siegel modular form}
A holomorphic function $f: \mathfrak{h} \to \mathbb{C}$ is a modular form of weight $k$ if:
$$f\left(\frac{a\tau+b}{c\tau+d}\right) = (c\tau+d)^k f(\tau)$$
for all $\begin{pmatrix} a & b \\ c & d \end{pmatrix} \in SL_2(\mathbb{Z})$.
\end{definition}

Key examples:
\begin{itemize}
\item Eisenstein series ($k \geq 2$): $E_{2k}(\tau) = 1 - \frac{4k}{B_{2k}}\sum_{n=1}^\infty \sigma_{2k-1}(n)q^n$ (weight-$2k$ modular forms)
\item \textbf{Quasi-modular Eisenstein series} $E_2$: $E_2(\tau) = 1 - 24\sum_{n=1}^\infty \sigma_1(n)q^n$ transforms as:
$$E_2(\gamma\tau) = (c\tau+d)^2 E_2(\tau) + \frac{6c(c\tau+d)}{\pi i}$$
for $\gamma = \bigl(\begin{smallmatrix} a & b \\ c & d \end{smallmatrix}\bigr) \in SL_2(\mathbb{Z})$.  The inhomogeneous correction term $6c(c\tau+d)/(\pi i)$ makes $E_2$ \emph{quasi-modular} (weight~$2$, depth~$1$); it appears throughout the genus-$1$ bar complex.
\item Dedekind eta: $\eta(\tau) = q^{1/24}\prod_{n=1}^\infty(1-q^n)$; key identity: $\vartheta_1'(0|\tau) = -2\pi\eta(\tau)^3$
\item Discriminant: $\Delta(\tau) = \eta(\tau)^{24} = q\prod_{n=1}^\infty(1-q^n)^{24}$
\end{itemize}

For genus $g$, Siegel modular forms are functions on the Siegel upper half-space $\mathfrak{h}_g$ transforming under $Sp_{2g}(\mathbb{Z})$.

\subsection{Elliptic Polylogarithms}
\index{elliptic polylogarithm|textbf}

The elliptic polylogarithms generalize classical polylogarithms:
\[
\text{Li}_n(z;\tau) = \sum_{k=1}^\infty \frac{q^k}{k^n}\frac{1}{1-zq^k}
\]

\begin{remark}[Convention]
Multiple non-equivalent definitions of ``elliptic polylogarithm'' exist in the literature (Levin, Zagier, Brown--Levin).  The $q$-series above is the convention used in this monograph, chosen for compatibility with the genus-$1$ bar differential.  For comparison with Zagier's Kronecker--Eisenstein series or Brown--Levin's iterated integrals on the universal elliptic curve, appropriate conversion factors are required.
\end{remark}

These appear in the genus $g$ bar differentials (where $g$ enters through the summation range):
\[
d^{(g)}_{\text{ell}} = \sum_{n=2}^{2g} \text{Li}_n(e^{2\pi iz};\tau) \cdot \eta^{\otimes n}
\]
